\section{Related Work}
\subsection{CNN-Based HSI Classification}
Convolutional neural networks are widely used for HSI classification because they provide strong local inductive bias. Early deep models processed spectral vectors or local image patches with one-dimensional, two-dimensional, or three-dimensional convolutions. Three-dimensional CNNs jointly encode spectral and spatial neighborhoods, which is helpful for distinguishing classes with similar spectra. Hybrid models such as HybridSN reduce the computational burden by combining 3-D and 2-D operations. Residual and dense connections further improve training stability. Nevertheless, CNN-based models mainly rely on local receptive fields, and their ability to represent long-range spectral or contextual relationships is limited unless deeper networks or larger patches are used.

\subsection{Transformer-Based HSI Classification}
Transformer architectures introduce self-attention to model long-range dependencies. In HSI classification, attention can be applied along spectral bands, spatial tokens, or joint spectral-spatial tokens. These models are effective for global context modeling and have become a strong family of HSI classifiers. However, standard self-attention has quadratic complexity with respect to token length. HSI data can contain hundreds of bands and many spatial tokens, so Transformer models may require substantial memory and training data. Lightweight, axial, grouped, and hierarchical attention mechanisms reduce the cost, but split sensitivity and label scarcity remain important concerns.

\subsection{Graph Neural Networks}
Graph neural networks model relational structure by message passing. For HSI classification, graph nodes can represent pixels, superpixels, regions, or learned tokens. Graph methods are useful because land-cover classes usually have spatial continuity and non-local similarity. Superpixel graphs can reduce noise and align better with homogeneous regions, while pixel graphs can preserve fine boundaries. However, graph construction can be sensitive to noise, and many graph models still require complementary spectral sequence modeling to capture detailed band dependencies.

\subsection{State-Space Models and Mamba}
State-space models provide a promising alternative to attention for long sequence processing. Mamba-style selective state-space models have attracted attention because they can capture long-range dependencies with linear complexity. In HSI classification, the ordered spectral bands naturally form a sequence, making state-space modeling a suitable mechanism for spectral representation. Recent Mamba-based HSI methods show the potential of this direction, but a direct replacement of attention with a state-space block is not sufficient for a strong journal contribution. The proposed DR-GSMamba extends the idea with graph reasoning, prototype uncertainty, and distributionally robust optimization.

\subsection{Robustness and Reproducibility}
Robustness is especially important in HSI classification because labeled samples are scarce and class distributions are imbalanced. Many papers report only one train-test split, which makes it difficult to judge whether improvements are statistically reliable. A robust HSI classifier should maintain performance under multiple random seeds, few-shot training ratios, label noise, and class imbalance. This paper therefore frames the proposed method around accuracy, stability, rare-class Macro-F1, and uncertainty rather than only a single overall accuracy number.
